\section{Introduction}
\label{sec:intro}

Minecraft, as the world's best-selling game, offers a range of tasks from exploration, survival to creating. It has become an important environment for researching efficient Reinforcement Learning (RL)~\cite{johnson2016malmo,guss2019minerl}. In particular, its extreme sparsity of rewards and huge complexity of the decision space pose significant challenges for RL. Currently, the most effective learning strategy involves pre-training through behavior cloning~\cite{baker2022video}, using learned behavioral priors to narrow the decision space. Nevertheless, it still requires billions of environmental interactions for effective learning due to the sparse reward nature of Minecraft.

On the other hand, previous researchers have proposed a variety of dense reward signals to enable efficient learning for specific sparse reward tasks~\cite{pathak2017curiosity,li2020random,oh2018self,andrychowicz2017hindsight,levy2017learning}. However, their applicability on the complex and long-horizon tasks in Minecraft remains an open question. To deeply reveal the challenges in Minecraft, we examine on several representative challenging tasks, \eg exploring underground for diamonds. We find that even after behavior cloning, most of these methods still fail to make significant progress on these tasks, further highlighting the difficulty of Minecraft and the limitations of existing dense reward methods.



It is noteworthy that for human players, Minecraft is a relatively simple casual game~\cite{duncan2011minecraft}. The advantage of human lies in their ability to summarize based on practice. For example, an accidental burning death from lava can teach human to avoid getting too close to it. Such summaries, based on life experience and practice, are key to human intelligence~\cite{tierney2009brain,seo2021impact}. Most existing RL methods overlook this ability. 
On the other hand, Large Language Models (LLMs) have recently demonstrated human-like common sense and reasoning capabilities~\cite{openai2023gpt,zhao2023survey,huang2022towards}. We find that leveraging LLMs can help RL agents simulate the practice summarization abilities of human. Based on the historical action trajectories and success-failure signals of the agents, LLMs can automatically design and refine corresponding auxiliary rewards, effectively overcoming the sparse reward challenge in Minecraft.

According to above analysis, we propose an automated method named Auto MC-Reward, to design and improve auxiliary reward functions according to task descriptions and historical action trajectories.
This method utilizes the task understanding and experience summarization abilities of LLMs, providing detailed and immediate rewards for learning guidance. Specifically
We first use LLMs to design task-related dense reward functions based on basic descriptions of the environment and tasks, named as Reward Designer. These reward functions are used to train agents after self-verification, \emph{i.e.} Reward Critic. To address potential biases or oversights in LLM's understanding, we also propose a LLM-based Trajectory Analyzer to analyze and summarize collected trajectories from the trained agent, and help Reward Designer to improve the reward functions. 

We verify the effectiveness of Auto MC-Reward on a series of representative benchmarks, including horizontal exploration for diamonds underground and approaching trees and animals in the plains biome. Experiments show that Auto MC-Reward achieves significantly better results on these tasks compared to original sparse reward and existing dense reward methods, showing its advanced ability of empowering efficient learning on sparse reward tasks. By iteratively refining the design of rewards functions, Auto MC-Reward enables the agent to efficiently learn new behaviors that is beneficial to the corresponding tasks, \emph{e.g.} avoiding lava, which greatly improves the success rate. Moreover, Auto MC-Reward achieves a high diamond obtaining success rate (36.5\%) with only raw information, demonstrating its ability of solving long-horizon tasks.
